\PassOptionsToPackage{quiet}{xeCJK}
\documentclass[withoutpreface,bwprint]{format}
\usepackage{etoolbox}
\BeforeBeginEnvironment{tabular}{\zihao{-5}}
% 参考文献改用纯tex格式，不再需要natbib+bibtex
\usepackage[framemethod=TikZ]{mdframed}
\usepackage{url}
\usepackage{array}
\usepackage{tabularx}
\usepackage{longtable}
\usepackage{fvextra}
\usepackage{placeins}
\usepackage{needspace}
\newcolumntype{C}{>{\centering\arraybackslash}X}
\newcolumntype{R}{>{\raggedleft\arraybackslash}X}
\newcolumntype{L}{>{\raggedright\arraybackslash}X}

\renewcommand{\textbf}[1]{{\song\bfseries #1}}

\usepackage{tikz}
\usetikzlibrary{arrows.meta}
\tikzset{
  box/.style={rectangle, draw, minimum width=3.2cm, minimum height=0.9cm, align=center},
  arrow/.style={thick, -{Stealth}}
}









\title{论文标题}



%%%%%%%%%%%%%%%%%%%%%%%%%%%%%%%%%%%%%%%%%%%%%%%%%%%%%%%%%%%%%
%% 正文
\begin{document}

\setcounter{page}{1}
\pagestyle{plain}

% 标题
{\centering\zihao{3}\song\bfseries 论文标题\par}

%%%%%%%%%%%%%%%%%%%%%%%%%%%%%%%%%%%%%%%%%%%%%%%%%%%%%%%%%%%%%
%% 章节骨架说明（示例结构，非官方要求）
%% 官方未规定正文章节划分；本骨架仅为起点，必须按题重组：
%% 1. 各章节可按需合并或删除——检验内容可并入各问题求解节，
%%    无推广价值时删除"8.模型改进推广"，禁止为凑齐骨架保留空章节；
%% 2. "5.模型的建立与求解.tex" 按实际问题数量重组为问题一/二/三……，
%%    不固定 5.1.x 子文件划分；
%% 3. 章节取舍依据见 references/paper-guides/cumcm-guide.md"推荐结构"一节。
%%%%%%%%%%%%%%%%%%%%%%%%%%%%%%%%%%%%%%%%%%%%%%%%%%%%%%%%%%%%%

\begin{abstract}

% ==== 开头段（1-2句话，直接切入题目任务，不复述背景意义）====
针对[题目核心问题]，本文基于[数据/机理]建立[模型体系概览]，依次求解[问题一/二/三/...]。

% ==== 逐问摘要（每问必须有方法和具体数值结果）====
\textbf{对于问题一,} 首先对...进行了...分析，接着采用...方法构建了...模型，最后通过...得到...。结果表明，...（关键数值，含单位）。

\textbf{对于问题二,} 首先...，接着...，最后...。结果表明...（具体数值）。

\textbf{对于问题三,} 首先...，接着...，最后...。结果表明...（具体数值）。

\textbf{对于问题四,} 首先...，接着...，最后...。结果表明...（具体数值）。

% ==== 结尾段（2句话，用检验结论收尾，禁止"性能优越/意义重大/决策参考"等空泛评价）====
通过[灵敏度分析/误差分析/蒙特卡洛模拟等]验证：[关键参数]在[范围]内变动时，[核心结果]保持稳定（[量化波动幅度]）。[一句话说明结果对题目场景的具体含义]。


\keywords{关键词1；关键词2；关键词3；关键词4；关键词5}

\end{abstract}



\section{引言}
\subsection{问题背景}

% 第1段：宏观背景 → 第2段：问题来源与困境 → 第3段：解决渴求


\subsection{问题重述}

\textbf{问题1：}...

\textbf{问题2：}...

\textbf{问题3：}...

\textbf{问题4：}...



\section{总体分析}

% 总述问题递进逻辑（纯文字，不画流程图）



%%%%%%%%%%%%%%%%%%%%%%%%%%%%%%%%%%%%%%%%%%%%%%%%%%%%%%%%%%%%%
\section{模型假设}

为简化问题，本文做出以下假设：

\begin{itemize}[itemindent=2em]
\item \textbf{假设1：}...
\item \textbf{假设2：}...
\item \textbf{假设3：}...
\item \textbf{假设4：}...
\end{itemize}



%%%%%%%%%%%%%%%%%%%%%%%%%%%%%%%%%%%%%%%%%%%%%%%%%%%%%%%%%%%%%
\section{符号说明}

本文使用的主要符号及其含义如表\ref{tab:符号说明}所示。

% longtable：可跨页，>{\centering\arraybackslash}p{} 居中，比例按内容调整（2列表总和0.92）
\begin{longtable}{>{\centering\arraybackslash}p{0.15\textwidth}>{\centering\arraybackslash}p{0.77\textwidth}}
  \caption{符号说明}
  \label{tab:符号说明} \\
  \toprule
  符号 & 说明 \\
  \midrule
  \endfirsthead
  \caption{符号说明（续）} \\
  \toprule
  符号 & 说明 \\
  \midrule
  \endhead
  \bottomrule
  \endfoot
  ... & ... \\
  ... & ... \\
  ... & ...
\end{longtable}



%%%%%%%%%%%%%%%%%%%%%%%%%%%%%%%%%%%%%%%%%%%%%%%%%%%%%%%%%%%%%
\section{模型的建立与求解}


%%%%%%%%%%%%%%%%%%%%%%%%%%%%%%%%%%%%%%%%%%%%%%%%%%%%%%%%%%%%%
\subsection{问题一的模型的建立和求解}


%%%%%%%%%%%%%%%%%%%%%%%%%%%%%%%%%%%%%%%%%%%%%%%%%%%%%%%%%%%%%
\subsubsection{具体分析}

% ==== 第1段：问题类型 + 大致方法 ====
问题一要求...，属于<评价/分类/预测/优化/机理分析>类问题。
本问采用...方法进行求解。

% ==== 第2段：完整步骤串联（首先→接着→然后→最后，不换段）====
首先对...进行...，
接着利用...方法...，
然后构建...模型...，
最后通过...得到...。

% ==== 第3段：引出流程图 ====
具体分析流程如图\ref{fig:analysis1_flow}所示。

% ==== 流程图（只改{}中文字，每框≤7字）====
\begin{figure}[H]
  \begin{center}
    \begin{tikzpicture}[x=1cm, y=0.7cm]
      \node[box] (n11) at (0, 0) {步骤1};
      \node[box] (n12) at (5, 0) {步骤2};
      \node[box] (n13) at (10, 0) {步骤3};
      \node[box] (n22) at (5, -3) {步骤5};
      \node[box] (n21) at (0, -3) {步骤4};
      \node[box] (n23) at (10, -3) {步骤6};
      \node[box] (n31) at (0, -6) {步骤7};
      \node[box] (n32) at (5, -6) {步骤8};
      \node[box] (n33) at (10, -6) {步骤9};
      \draw[arrow] (n11) -- (n12);
      \draw[arrow] (n12) -- (n13);
      \draw[arrow] (n23) -- (n22);
      \draw[arrow] (n22) -- (n21);
      \draw[arrow] (n13) -- ++(0,-1.0) -| (n23);
      \draw[arrow] (n21) -- (n31);
      \draw[arrow] (n31) -- (n32);
      \draw[arrow] (n32) -- (n33);
    \end{tikzpicture}
    \caption{问题一分析流程图}
    \label{fig:analysis1_flow}
  \end{center}
\end{figure}

% ============================================================
\subsubsection{模型准备}

% ==== 引文：简述本节做什么 ====
在进行建模之前，需要...。本节将从数据预处理和算法选择两方面进行准备。

% ==== 数据预处理（仅问题1写此段）====
原始数据来源于...，包含...条记录和...个字段。
经检查发现以下问题：...（缺失值/量纲不统一/异常值/编码问题）。

针对上述问题，按以下步骤进行预处理：
第一步，...（如缺失值处理，具体方法+参数）。
第二步，...（如标准化/归一化，写出公式）。
第三步，...（如特征编码/衍生）。

预处理后，数据规模为...行×...列，各指标统计量如下表所示。

% 预处理对比表（可选）
\begin{table}[H]
  \centering
  \caption{数据预处理前后统计量对比}
  \begin{tabularx}{\textwidth}{CCCCC}
    \toprule
    指标 & 预处理前均值 & 预处理前标准差 & 预处理后均值 & 预处理后标准差 \\
    \midrule
    ... & ... & ... & ... & ... \\
    \bottomrule
  \end{tabularx}
  \label{tab:预处理对比}
\end{table}

综上所述，数据预处理完成，为后续建模提供了清洁的输入数据。

% ==== 算法选择 ====
该问题本质上是一个...问题，需要在...条件下...。
常用求解算法包括A算法、B算法、C算法等。本节从维度1、维度2、维度3三方面进行评估。

\begin{table}[H]
  \centering
  \caption{算法能力对比}
  \begin{tabularx}{\textwidth}{CCCCC}
    \toprule
    算法名称 & 维度1 & 维度2 & 维度3 & 稳定性 \\
    \midrule
    算法A（缩写） & 优 & 良 & 中 & 良 \\
    算法B（缩写） & 良 & 优 & 良 & 优 \\
    算法C（缩写） & 中 & 中 & 优 & 中 \\
    \bottomrule
  \end{tabularx}
  \label{tab:算法对比1}
\end{table}

通过对比可以看出，算法A在维度1方面表现最优，...。
因此本文选择算法A作为核心求解方法。

综上所述，本节完成了数据预处理和算法选型，接下来将建立数学模型并进行求解。



%%%%%%%%%%%%%%%%%%%%%%%%%%%%%%%%%%%%%%%%%%%%%%%%%%%%%%%%%%%%%
\subsubsection{模型建立}

% ==== 开篇引文 ====
针对问题一，本节将从...角度建立...模型。具体过程如下。

% ==== 步骤1：<标题> ====
\textbf{步骤1：...}

首先，...。设...为...，则有：
\begin{equation}
    ... = ...
\end{equation}
其中，$...$表示...。

将...代入...，可得：
\begin{equation}
    ... = ...
\end{equation}
进一步展开得：
\begin{equation}
    ... = ...
\end{equation}

% ==== 步骤2：<标题> ====
\textbf{步骤2：...}

根据...的性质，可以建立以下关系：
\begin{equation}
    ... = ...
\end{equation}
其中，$...$为...，$...$为...。

由式(1)和式(2)联立可得：
\begin{equation}
    ... = ...
\end{equation}

% ==== 步骤3：<标题> ====
\textbf{步骤3：...}

% ...（按需增加步骤，每步≥6个公式）...

% ==== 步骤4：求解算法 ====
\textbf{步骤4：...算法}

% 算法流程公式 + 参数表

\begin{table}[H]
  \centering
  \caption{算法参数设置}
  \begin{tabularx}{\textwidth}{CCCC}
    \toprule
    参数名 & 参数含义 & 取值 & 选取依据 \\
    \midrule
    ... & ... & ... & ... \\
    \bottomrule
  \end{tabularx}
  \label{tab:参数1}
\end{table}

% ==== 总结桥接 ====
综上所述，本节完成了...模型的建立，得到了...（目标函数/评价体系/预测模型），接下来将基于此进行求解。

% ============================================================
\subsubsection{模型求解}

% ==== 引文 ====
基于上述模型，本节利用...数据进行求解。求解环境为...，核心参数如表\ref{tab:参数1}所示。

% ==== 结果表格 ====
求解得到以下结果：

\begin{table}[H]
  \centering
  \caption{...结果}
  \begin{tabularx}{\textwidth}{CCCC}
    \toprule
    指标1 & 指标2 & 指标3 & 指标4 \\
    \midrule
    ... & ... & ... & ... \\
    \bottomrule
  \end{tabularx}
  \label{tab:结果1}
\end{table}

% ==== 结果图（按需启用下方示例，并在正文增加对应引用）====

% 图引用示例：
% \begin{figure}[H]
%   \centering
%   \includegraphics[width=0.8\textwidth]{../求解/问题一/figures/xxx.png}
%   \caption{...}
%   \label{fig:xxx}
% \end{figure}

% ==== 结果分析 ====
由表\ref{tab:结果1}可知，...。若加入结果图，应在此解释图中的具体趋势与数值。
综合以上结果，...（一段分析总结）。



% 如有更多问题，按以下格式新增：
% \input{5.2.问题2的建立求解.tex}
% \input{5.3.问题3的建立求解.tex}



%%%%%%%%%%%%%%%%%%%%%%%%%%%%%%%%%%%%%%%%%%%%%%%%%%%%%%%%%%%%%
\section{模型检验与分析}

为验证本文构建的模型的有效性和稳健性，进行以下检验。

\subsection{误差分析}

\begin{table}[H]
  \centering
  \caption{模型误差分析结果}
  \label{tab:error}
  \begin{tabularx}{\textwidth}{CCCC}
    \toprule
    指标 & 训练集 & 测试集 & 说明 \\
    \midrule
    $R^2$ & ... & ... & 决定系数 \\
    RMSE & ... & ... & 均方根误差 \\
    MAE & ... & ... & 平均绝对误差 \\
    MAPE(\%) & ... & ... & 平均绝对百分比误差 \\
    \bottomrule
  \end{tabularx}
\end{table}

如表\ref{tab:error}所示，...

\subsection{灵敏度分析}

\begin{table}[H]
  \centering
  \caption{灵敏度分析结果}
  \label{tab:sensitivity}
  \begin{tabularx}{\textwidth}{CCC}
    \toprule
    参数变化 & 输出变化 & 灵敏度评估 \\
    \midrule
    ... & ... & ... \\
    \bottomrule
  \end{tabularx}
\end{table}

如表\ref{tab:sensitivity}所示，...



%%%%%%%%%%%%%%%%%%%%%%%%%%%%%%%%%%%%%%%%%%%%%%%%%%%%%%%%%%%%%
\section{模型的评价}

\subsection{模型的优点}
\begin{itemize}[itemindent=2em]
\item \textbf{优点1：}
\item \textbf{优点2：}
\item \textbf{优点3：}
\end{itemize}

\subsection{模型的缺点}
\begin{itemize}[itemindent=2em]
\item \textbf{缺点1：}
\item \textbf{缺点2：}
\end{itemize}



%%%%%%%%%%%%%%%%%%%%%%%%%%%%%%%%%%%%%%%%%%%%%%%%%%%%%%%%%%%%%
\section{模型的改进与推广}

\subsection{模型的改进}

...

\subsection{模型的推广}

...



%%%%%%%%%%%%%%%%%%%%%%%%%%%%%%%%%%%%%%%%%%%%%%%%%%%%%%%%%%%%%
%% 参考文献

\begin{thebibliography}{99}

% 必须根据论文实际内容生成，禁止保留模板示例文献
% 每条格式：\bibitem{标签} 作者. 题名. 出版信息.
% 标签：作者姓+年份+关键词，如 \bibitem{behzadian2012topsis}
% 数量8-15条，每条必须在正文有\cite{}引用，无引用的删掉

\bibitem{ref1} ...
\bibitem{ref2} ...
\bibitem{ref3} ...

\end{thebibliography}




\makeatletter
\let\@oldaddcontentsline\addcontentsline
\renewcommand{\addcontentsline}[3]{}
\section*{附录}
\let\addcontentsline\@oldaddcontentsline
\makeatother

\begin{table}[H]
  \centering
  \caption{附件说明表}
  \label{tab:appendix}
  \renewcommand{\arraystretch}{1.2}
  \begin{tabularx}{\textwidth}{LX}
    \toprule
    文件 & 说明 \\
    \midrule
    问题一\_xxx.py & 问题1求解代码 \\
    xxx.png & xxx图 \\
    \bottomrule
  \end{tabularx}
\end{table}

\subsection*{附录2 源程序代码}

%（官方格式规范第五条要求）此处粘贴本次建模用到的全部完整、可运行的源程序代码。
%（每个问题的求解脚本从 import 到输出完整粘贴，含 mrite_runtime.py；若确实没有用到
%程序，改为注明“本论文没有用到程序”）。
%支撑材料（含源程序、自主查阅的数据资料、较大篇幅中间结果图表）压缩为RAR/ZIP
%单独上传，大小不超过20MB，文件列表已在上表列出。



\end{document}
